% =====================================================================
% Appendix: Measurement uncertainty of the steady-state results
% Calculation: figures/data/make_rcp_data.py (also writes the plotted data)
% =====================================================================
\chapter{Measurement Uncertainty}
\label{app:uncertainty}

This appendix derives the uncertainties quoted in \cref{tab:measured-performance}
and shown as error bars in \cref{fig:galinstan_rcp,fig:water_rcp,fig:galinstan_rcp_cfd}.
The instrument accuracies are propagated through the thermal network of
\cref{eq:network} to the heat reaching the plate, the total resistance and the
inferred cold-plate resistance of \cref{eq:rcp-inferred}. The calculation is
scripted in \texttt{figures/data/make\_rcp\_data.py}, which also writes the data
files plotted in those figures.

\section{Instrument Accuracy}

\Cref{tab:app-instruments} lists the accuracy assumed for each measured
quantity. The thermocouple accuracy dominates, because every resistance is
obtained from a temperature difference of only \SIrange{16}{30}{\celsius}.

\begin{table}[htbp]
  \centering
  \caption{Accuracy of the measured quantities.}
  \label{tab:app-instruments}
  \small
  \begin{tabularx}{\linewidth}{@{}>{\raggedright\arraybackslash}p{3.3cm}
      >{\raggedright\arraybackslash}p{4.2cm} L@{}}
    \toprule
    Quantity & Instrument & Uncertainty \\
    \midrule
    Steady-state temperature & K-type thermocouples
      & $\pm\SI{2.2}{\celsius}$ per reading~\cite{loh_project_2025} \\
    Heater power, $\Qin$ & GW Instek GPE-3323 readback
      & $\pm\SI{2}{\percent}$ (assumed) \\
    Resistance to ambient, $\Rmisc$ & No-flow tests at \SI{12}{\watt}
      & $\pm\SI{8}{\percent}$ (Galinstan), $\pm\SI{9}{\percent}$ (water),
        propagated \\
    Galinstan flow & Energy balance; EMF sensor as a check
      & $\pm\SI{18}{\percent}$ (energy balance, propagated) \\
    Water flow & Sea Zhongjiang digital flow meter
      & Not specified; out of range below \SI{1}{\litre\per\minute} \\
    AC resistance & Agilent E4980A LCR meter
      & $\pm\SI{1.5}{\milli\ohm}$ (repeatability band) \\
    \bottomrule
  \end{tabularx}
\end{table}

\section{Propagation}

The errors of the individual readings are treated as independent and combined
by root-sum-square, which is a conservative choice for a temperature
difference (\cref{sec:app-unc-interpretation}). Any difference between two
thermocouple readings therefore has an uncertainty of
\begin{equation}
  u(\Delta T) = \sqrt{2}\,u_T = \sqrt{2} \times \SI{2.2}{\celsius}
    = \SI{3.1}{\celsius}.
\end{equation}
The resistance to ambient, $\Rmisc = (\Tres - T_{\mathrm{A}})/Q_0$, was measured
at $Q_0 = \SI{12}{\watt}$ with a temperature rise of about \SI{40}{\celsius}
(Galinstan) and \SI{35}{\celsius} (water), giving the relative uncertainties in
\cref{tab:app-instruments}. The heat lost to ambient during a test,
$Q_{\mathrm{misc}} = (\Tres - T_{\mathrm{A}})/\Rmisc$, and the heat reaching the
plate then have uncertainties
\begin{equation}
  \frac{u(Q_{\mathrm{misc}})}{Q_{\mathrm{misc}}}
    = \sqrt{\left(\frac{u(\Delta T)}{\Tres - T_{\mathrm{A}}}\right)^{2}
      + \left(\frac{u(\Rmisc)}{\Rmisc}\right)^{2}},
  \qquad
  u(\Qplate) = \sqrt{u(\Qin)^{2} + u(Q_{\mathrm{misc}})^{2}}.
\end{equation}
The total resistance $R = (\Tres - \Tin)/\Qplate$ and the inferred cold-plate
resistance $\Rcp = R - \Rpack$ have the same absolute uncertainty,
\begin{equation}
  u(\Rcp) = u(R) = R\,\sqrt{\left(\frac{u(\Delta T)}{\Tres - \Tin}\right)^{2}
      + \left(\frac{u(\Qplate)}{\Qplate}\right)^{2}},
  \label{eq:app-u-r}
\end{equation}
with $\Rpack$ treated as exact. Its own calibration uncertainty, about
$\pm\SI{0.05}{\celsius\per\watt}$ from the same expression, would shift every
inferred $\Rcp$ by the same amount and does not change the comparison between
tests. The Galinstan flow from the energy balance,
$\Vdot = \Qplate/[\rho\cp(\Tout - \Tin)]$, has a relative uncertainty of
$\sqrt{(u(\Qplate)/\Qplate)^{2} + (u(\Delta T)/(\Tout - \Tin))^{2}}$, about
\SI{18}{\percent} for a \SIrange{17}{18}{\celsius} coolant temperature rise. The
water flow meter has no stated accuracy, so no flow uncertainty is propagated for
water.

\section{Results}

\Cref{tab:app-uncertainty} gives the propagated uncertainty for each
steady-state test. The total resistance is known to about \SI{11}{\percent} for
full-flow Galinstan, \SI{14}{\percent} for low-flow Galinstan and
\SI{20}{\percent} for water, so the difference between the two coolants is far
larger than its uncertainty. For water, however, the uncertainty of
$\pm\SI{0.05}{\celsius\per\watt}$ exceeds the inferred $\Rcp$ of
\SIrange{0.015}{0.025}{\celsius\per\watt}, because $\Rcp$ is a small difference
between two resistances of about \SI{0.25}{\celsius\per\watt}.

\begin{table}[htbp]
  \centering
  \caption{Propagated uncertainty of the steady-state results.}
  \label{tab:app-uncertainty}
  \small
  \begin{tabular}{@{}lcccc@{}}
    \toprule
    Coolant & $\Vdot$ & $\Qplate$ & $(\Tres - \Tin)/\Qplate$ & $\Rcp$ \\
     & (\si{\milli\litre\per\second}) & (\si{\watt}) & (\si{\celsius\per\watt})
      & (\si{\celsius\per\watt}) \\
    \midrule
    Galinstan & $1.45 \pm 0.26$ & $60.4 \pm 2.0$ & $0.488 \pm 0.054$ & $0.246 \pm 0.054$ \\
    Galinstan & $0.77 \pm 0.15$ & $30.8 \pm 1.4$ & $0.799 \pm 0.108$ & $0.557 \pm 0.108$ \\
    Water     & 58.3            & $62.6 \pm 2.0$ & $0.257 \pm 0.050$ & $0.015 \pm 0.050$ \\
    Water     & 29.2            & $62.8 \pm 2.0$ & $0.264 \pm 0.050$ & $0.022 \pm 0.050$ \\
    Water     & 12.2\tnote{a}   & $62.3 \pm 2.0$ & $0.264 \pm 0.051$ & $0.022 \pm 0.051$ \\
    Water     & 8.3\tnote{a}    & $62.2 \pm 2.0$ & $0.267 \pm 0.051$ & $0.025 \pm 0.051$ \\
    \bottomrule
  \end{tabular}
  \floatnote{Galinstan flow from the energy balance. \tnote{a}Below the flow
    meter's range.}
\end{table}

\section{Galinstan Flow Sensor}

The EMF flow sensor gives an independent check of the Galinstan flow.
\Cref{tab:app-emf} lists its readings against pump current, converted with
$\Vdot = Eh/B$~\cite{kwon_project_2025}, together with the pump power measured
across its terminals. At \SI{18}{\ampere} and \SI{9}{\ampere} the sensor gives
1.37 and \SI{0.74}{\milli\litre\per\second}, which the report rounds to 1.4 and
\SI{0.75}{\milli\litre\per\second}. The energy balance gives 1.45 and
\SI{0.77}{\milli\litre\per\second}, \SIrange{4}{6}{\percent} higher, well within
its own uncertainty. The flow rises almost linearly with current, as
\cref{eq:mhd-pressure} predicts for a pump working against a laminar loop.

The conversion from EMF to flow assumes the flux density $B$ of the sensor
magnets, so the sensor is only accurate for the magnets it was calibrated with.
The magnets were replaced during the project, which changed the EMF at the same
pump current by about \SI{10}{\percent} (from 135.5 to \SI{149}{\micro\volt} at
\SI{18}{\ampere}), and the project's calculation sheets assume between
\SI{0.38}{\tesla} and \SI{0.45}{\tesla}. If the magnets are swapped without
recalibrating the sensor, the measured flow can therefore be in error by up to
about \SI{18}{\percent}.

\begin{table}[htbp]
  \centering
  \caption{EMF flow-sensor readings and pump power against pump current.}
  \label{tab:app-emf}
  \small
  \begin{tabular}{@{}cccc@{}}
    \toprule
    Pump current & Sensor EMF & Flow & Pump power \\
    (\si{\ampere}) & (\si{\micro\volt}) & (\si{\milli\litre\per\second}) & (\si{\watt}) \\
    \midrule
    9.0  & 80.5  & 0.74 & 0.21 \\
    12.0 & 103.3 & 0.95 & 0.37 \\
    18.0 & 149.0 & 1.37 & 0.83 \\
    18.2 & 150.5 & 1.39 & 0.85 \\
    \bottomrule
  \end{tabular}
    \floatnote{Room temperature, replacement sensor magnets, $B = \SI{0.38}{\tesla}$
    assumed. Source: project calculation sheet.}
\end{table}

\section{Interpretation}
\label{sec:app-unc-interpretation}

The $\pm\SI{2.2}{\celsius}$ accuracy of a thermocouple is mainly a fixed
calibration offset rather than random noise. Treating the readings as
independent therefore gives an upper bound on the uncertainty of any single
result. Because the same resistors and thermocouples, and the same calibrated
$\Rpack$, are used in every test, much of this offset is common to all tests, so
differences between tests are resolved better than the error bars suggest. The
water results support this: the measured resistor temperature rise varied by only
\SI{0.5}{\celsius} (\SIrange{16.1}{16.6}{\celsius}) across a seven-fold change in
flow. The repeatability of the tests themselves could not be assessed, because no
test was repeated (\cref{sec:disc-limitations}). The CHT results have a numerical
uncertainty of about \SI{0.25}{\celsius} from the mesh
(\cref{app:cht_validation}), or about $\pm\SI{0.004}{\celsius\per\watt}$ in
$\Rcp$, which is smaller than the plotted markers.
